%==============================================================================
% CHAPTER 3
%==============================================================================
\chapter{Towards a fuel-rod digital twin}
\markboth{Chapter 3 : Towards a fuel-rod digital twin}{}

\section{Scoping : what this twin is, and what it is not}
\label{sec:cadrage}

We should say first what we mean by a ``digital twin of a power plant'', because
the phrase promises more than what we built. A complete digital twin of a plant
would handle the full coupled
multiphysics : neutronics, thermal-hydraulics and fuel performance. It would
thereby be able to predict the state of a fuel rod from the state of the core,
and conversely. OFFBEAT, however, only deals with fuel performance. What we built is therefore the digital twin
of a \textbf{single fuel rod}, and everything outside that domain enters as a
boundary condition.

The second limitation concerns ``real time''. OFFBEAT is a batch solver, one run
of which takes from a few seconds to several hours, and there is no sensor
measuring the state of a rod inside the core in real time. In practice there is
no real time but rather an \textbf{update-driven} refresh of the state : at each
new operating history, the system re-simulates and updates its prognosis. True
real time can only come from a \textbf{reduced-order model}, hence block~D5.

Table~\ref{tab:agent-vs-twin} summarises the shift.

\begin{table}[htbp]
\centering
\caption[Base agent versus digital-twin prototype]{Differences between the base
agent and the digital-twin prototype aimed at here.}
\label{tab:agent-vs-twin}
\begin{tabularx}{0.95\textwidth}{@{}lXX@{}}
\toprule
 & \textbf{Base agent} & \textbf{Digital twin}\\
\midrule
\textbf{When} & after a completed simulation & at each new piece of data\\
\textbf{What} & creates, runs, post-processes & + monitors safety thresholds\\
\textbf{Behaviour} & reactive (repairs a crash that occurred) &
  predictive (announces a hazard before it occurs)\\
\bottomrule
\end{tabularx}
\end{table}

\section{D1 : The safety analyser}

The first block reads a completed simulation and assigns, criterion by
criterion, a status : safe, close to the threshold, or exceeded. It follows
the same pattern as self-healing : an editable JSON knowledge base, a
deterministic treatment, and an optional language model fallback to explain in
plain language the criteria that are exceeded.

Each criterion describes the field to read, the reduction to apply (maximum,
minimum, largest-magnitude value), the mesh region concerned, the limit and its
direction, that is whether the quantity must stay below or above the threshold.
This last point
proved trickier than it appears : gap width is the only criterion whose danger is
a value that is \textit{too small}, and this special case was the origin of a
defect detected late (§\ref{sec:bugs}).

The tresholds written in the base is an indicative value as we were not able to find
a better criteria. They are published design criteria and not purely invented but
that would need more expertise to be more accurate.


\section{D2 : Prognosis}

The second block exploits the fact that a simulation writes several time steps.
For each non-green criterion, we extract the time evolution of the peak, fit a
linear trend to it, and read off the instant at which the threshold would be
extrapolated.

The cautious rule adopted is to announce nothing when the trend is not moving
towards the limit : if the quantity is moving away from it, or is stationary, the
pronostic says so explicitly rather than producing a meaningless extrapolation.
This extrapolation is of first-order, hence valid as long as there is no regime change.

\section{D3 : Operating-history update}

The third block closes the loop : an operating history file is monitored, its
last line representing the current state. At each update, the case is rewritten
with the new parameters, the earlier time steps are cleaned, the simulation is
re-run and the safety analysis is replayed. This cycle is the "update-driven" part.

We first called this block \emph{data assimilation}, and that name was wrong.
Assimilation combines an observation with the model state and updates the
uncertainty on that state, as a Kalman filter does. Nothing of the sort happens
here : the operating history is an \textbf{input} that we replay, not a
measurement we reconcile with a prediction. Genuine assimilation would require
in-core measurements, which do not exist for a single rod, and an explicit
representation of the model uncertainty. It stays a longer-term goal.

\section{D4 : The dashboard}

The web interface was extended with a safety panel. There are per-criterion gauges,
refreshed periodically, showing the current value, the limit and the fraction of
the threshold reached. To this is added an interactive slider-based simulator,
described below.

\section{D5 : The surrogate}

\subsection{Motivation and principle}


The dashboard highlights a paradox about the analysis being instantaneous, but using an already runned simulation.
Then, exploring a new scenario such as a ten per cent increase in power requires re-running the
solver.

The aim of block D5 is to break through this limitation by providing a \textbf{reduced-order model} of the solver.
A design of experiments is built by crossing two
parameters, linear heat rate and irradiation duration. For each combination a
complete OFFBEAT simulation is run and three safety quantities are recorded. A
statistical model is then fitted to this dataset.

The choice of the sampling domain turned out to be the main decision of this
whole block, and we got it wrong. Our first version only swept durations of a few
thousand seconds, because a short design of experiments was cheap to build. But
the rod reaches thermal equilibrium after about \SI{5000}{\second} : past that
point, changing the duration changes almost nothing. The problem we were actually
fitting had only one variable left, and section~\ref{sec:validation-emulateur}
shows what that cost us.



\subsection{Choice of a Gaussian process}

The model adopted is a \textbf{Gaussian process} (kriging), with a Matérn
kernel. This choice, preferred over polynomial regression, is justified by two
properties well suited to the very small number of available points :

\begin{itemize}
  \item a Gaussian process \textbf{interpolates} the training points, which is
    consistent when they come from a deterministic code rather than from noisy
    measurements;
  \item it provides an \textbf{uncertainty} at every prediction point. This
    error bar is propagated all the way to the interface : the user sees not only
    the predicted value but the confidence that can be placed in it, which is
    indispensable as soon as a safety computation is replaced by an
    approximation.
\end{itemize}

We validate by \textit{leave-one-out} : each point is predicted by
a model refitted on all the others. This is an honest measure of generalisation
ability, far more meaningful than a goodness-of-fit coefficient over a few dozen
points. The results are presented in the next chapter, where it will be seen
that a statistically honest validation can nonetheless mislead if the sampled
domain is not itself physically relevant.

\clearpage

\clearpage
